% The double-stranded lattice sum as a fermionic partition function   (Bio 13 / XXVI)
% Build: pdflatex anccft26 (x3)
\documentclass[11pt]{article}
\usepackage{lmodern}
\usepackage[T1]{fontenc}
\usepackage{amsmath,amssymb,amsthm,mathtools}
\usepackage[margin=1.1in]{geometry}
\usepackage{booktabs}
\usepackage{microtype}
\usepackage[hidelinks]{hyperref}

\theoremstyle{plain}
\newtheorem{theorem}{Theorem}
\newtheorem{proposition}[theorem]{Proposition}
\newtheorem{lemma}[theorem]{Lemma}
\newtheorem{corollary}[theorem]{Corollary}
\theoremstyle{definition}
\newtheorem{definition}[theorem]{Definition}
\theoremstyle{remark}
\newtheorem{remark}[theorem]{Remark}

\newcommand{\Z}{\mathbb{Z}}
\newcommand{\rank}{\operatorname{rank}}
\newcommand{\spec}{\operatorname{spec}}
\newcommand{\supp}{\operatorname{supp}}
\newcommand{\Loc}{\operatorname{Loc}}
\newcommand{\Nseq}{N_{\mathrm{seq}}}
\newcommand{\DS}{\mathrm{DS}}
\newcommand{\Lm}{\Lambda^{-}}
\newcommand{\bpsi}{\bar\psi}

\title{\bfseries The double-stranded lattice sum\\ as a fermionic partition function}
\author{Ruqing Chen\\[2pt]
\normalsize GUT Geoservice Inc., Rivi\`ere-Beaudette, Qu\'ebec J0P 1R0, Canada\\
\normalsize\texttt{ruqing@hotmail.com}}
\date{22 September 2026}

\begin{document}
\maketitle

\begin{abstract}
\noindent
Bio~12 counted the lattice points of the double-stranded box of $\phi$X174 at $k=10$, and
those with connected support, but not the number $\#\DS$ of double-stranded
reconstructions. The obstacle is the arborescence factor $t(G_c)$, a determinant of the whole
support. We show that it is a fermionic partition function with one \emph{even} factor per
arc, so that it can be accumulated arc by arc in any order. A reduced determinant needs a
root, and the support moves with the point. Both difficulties are removed together by the
reverse-complement symmetry $c\mapsto m-c$, which fixes the summand and halves the box along
any orbit of odd multiplicity, where one fixed vertex is always in the support. The weight is
made integral by trading $\prod_a c(a)!^{-1}$ for binomials and one constant. Everything is
checked exactly against Bio~8 at $k=11,12$ and on random covers. Frontier elimination of the
fermionic weight is exact but does not reach $k=10$, because its state is linear and a single
history already fills up to $3^{w}$ of its coordinates. Nor does the positive (arborescence)
form of the same weight: it has one state per history, but merges too few of them. The number
is obtained instead by visiting all $1\,240\,843\,950$ points of the half-box in compiled
code, contracting forced arcs so that only a small determinant is left at each point, and
working modulo five primes. The run reproduces both of Bio~12's counts, halved, and agrees with
the exact arborescence elimination on $k=10$ sub-boxes. The result is
\[
\#\DS(\phi\mathrm{X}174,10)=31\,925\,753\,246\,212 ,
\]
the last blank entry of Bio~8's Table~2.
\end{abstract}

\medskip
\noindent\textbf{Keywords.} de Bruijn graph, reverse-complement double cover, circulation
lattice, matrix-tree theorem, Berezin integral, Grassmann variables, transfer matrix,
frontier elimination, Eulerian circuits, multimodular arithmetic, genome assembly.

\section{What is left to count}\label{sec:left}

We keep the notation of \cite{Bio8} and \cite{Bio12}. $S$ is a circular sequence, $D_k$ the
labelled reverse-complement double cover, $A$ its set of arc contents, $m(a)$ the number of
labelled arcs of content $a$, and $\rho$ the reverse-complement involution, so that
$m(\rho a)=m(a)$ and $D_k$ is balanced. A circular sequence $T$ is a \emph{double-stranded
reconstruction} of $S$ at word length $k$ when $\spec(T)+\spec(\bar T)=\spec(S)+\spec(\bar S)$;
writing $c=\spec(T)$ this says that $c$ is a circulation with $c+\rho c=m$, and with
$f=2c-m$ that $f$ lies in the $(-1)$-eigenlattice $\Lm=\{f\in Z_1(D_k):\rho f=-f\}$, inside the
box $|f|\le m$ and in the parity class of $m$. Hence \cite[Thm.~7]{Bio8}
\begin{equation}\label{eq:sum}
\#\DS(S,k)\;=\;\sum_{f}\Nseq\Big(\frac{m+f}{2}\Big),
\qquad
\Nseq(c)=\frac{t(G_c)\,\Loc(G_c)}{\prod_a c(a)!},
\end{equation}
where $G_c$ is the multidigraph with $c(a)$ copies of each content, $t(G_c)$ its number of
spanning arborescences to a root, and $\Loc(G_c)=\prod_v(d^+_c(v)-1)!$.

Bio~12 observed that the constraint defining the sum is one linear equation per vertex, each
involving only the four endpoints of one $\rho$-orbit, and replaced Bio~8's backtracking by a
\emph{frontier elimination}: process the orbits in a good order carrying, as state, the partial
imbalance at the vertices still open. At $k=10$, where $\rank\Lm=48$, that counts the
$2\,481\,687\,900$ points of the box in seconds with a frontier $26$ wide and $50\,758$ states,
and, carrying in addition the partition of the frontier into the classes already joined, the
$909\,103\,210$ points whose support is connected.

It did not reach $\#\DS$. The reason is entirely in $t(G_c)$: the number of spanning
arborescences is the determinant of a reduced Laplacian of the whole support --- of order $125$
at $k=10$ --- and a determinant is not a product of local scalars. A frontier state that is a
number, or a partition, cannot carry it.

The point of this note is that it \emph{is} a product of local \emph{fermions}.

\section{The determinant as a product over arcs}\label{sec:berezin}

Let $V$ be the vertex set, $r\in V$ a distinguished \emph{root}, and let
$\Lambda=\Lambda\big(\bpsi_v,\psi_v: v\in V\smallsetminus\{r\}\big)$ be the Grassmann algebra on
two odd generators per non-root vertex; we put $\psi_r=\bpsi_r=0$. For $c\in\Z_{\ge0}^A$ let
$L_c$ be the out-Laplacian of $G_c$, $L_c(v,v)=d^+_c(v)-(\text{loops at }v)$ and
$L_c(v,u)=-\sum_{a:\,t(a)=v,\,h(a)=u}c(a)$ for $u\ne v$, and let
\begin{equation}\label{eq:Aprime}
A_c\;=\;L_c\;+\;\operatorname{diag}\big(\mathbf 1_{v\notin\supp c}\big),
\qquad \supp c=\{v: d^+_c(v)>0\},
\end{equation}
so that $A_c$ agrees with $L_c$ on $\supp c$ and is the identity off it. Write
$\widetilde A_c^{\,r}$ for $A_c$ with the row and the column of $r$ deleted.

\begin{proposition}[arc factorisation]\label{prop:arc}
In $\Lambda$,
\begin{equation}\label{eq:prod}
P_c \;=\;
\prod_{a\in A}\Big(1-c(a)\,\bpsi_{t(a)}\big(\psi_{t(a)}-\psi_{h(a)}\big)\Big)
\cdot\!\!\prod_{v\notin\supp c}\!\!\big(1-\bpsi_v\psi_v\big)
\;=\;(-1)^{|V|-1}\,\det \widetilde A_c^{\,r}\cdot\Psi ,
\end{equation}
where $\Psi=\bpsi_{v_1}\psi_{v_1}\bpsi_{v_2}\psi_{v_2}\cdots$ runs over
$V\smallsetminus\{r\}$ in any order. Every factor in \eqref{eq:prod} has even degree, so the
product may be formed in any order whatsoever, one arc at a time.
\end{proposition}

\begin{proof}
Each factor is $1+x$ with $x$ a linear combination of $\bpsi_v\psi_u$, so $x^2=0$ and
$1+x=e^{x}$; and $\bpsi_v^2=0$ makes the product of the factors belonging to one tail $v$ equal
to $1-\bpsi_v\sum_u A_c(v,u)\psi_u$, the loops cancelling between the two terms of
$\psi_{t(a)}-\psi_{h(a)}$. Hence
$P_c=\prod_{v\ne r}\big(1-\bpsi_v(A_c\psi)_v\big)=\exp\big(-\sum_{v,u}\bpsi_vA_c(v,u)\psi_u\big)$,
the Berezin integrand of $\det\widetilde A^{\,r}_c$. Expanding,
$P_c=\sum_\sigma\prod_{v}\big(-A_c(v,\sigma v)\big)\,
\bpsi_{v_1}\psi_{\sigma v_1}\bpsi_{v_2}\psi_{\sigma v_2}\cdots$,
the sum over the bijections $\sigma$ of $V\smallsetminus\{r\}$. Each block $\bpsi_v\psi_{\sigma v}$
is even, so the blocks commute and the vertex order chosen for $\Psi$ is immaterial; and
interchanging two of the $\psi$'s moves an odd generator across a block of odd length, which
costs $-1$, so the monomial equals $\operatorname{sgn}(\sigma)\Psi$. The $(-1)^{|V|-1}$ is the
one minus sign per row.
\end{proof}

\begin{remark}\label{rem:interleave}
The constant $(-1)^{|V|-1}$ in \eqref{eq:prod} is attached to the \emph{interleaved} order
$\bpsi_{v_1}\psi_{v_1}\bpsi_{v_2}\psi_{v_2}\cdots$. Two interleaved orders differ by a
permutation of the even blocks $\bpsi_v\psi_v$ and so give the same constant, but listing all
the $\bpsi$ first and all the $\psi$ afterwards does not: that costs a shuffle sign which, when
the two lists are ordered differently --- as they are in an implementation where $\bpsi_v$ and
$\psi_v$ have different lifetimes --- varies from instance to instance. This is not a
scruple; it is the one place where an implementation of \eqref{eq:prod} silently returns
$-\#\DS$.
\end{remark}

The point of Proposition~\ref{prop:arc} for the frontier is this. The factor of the arc $a$
involves only $\bpsi_{t(a)}$, $\psi_{t(a)}$ and $\psi_{h(a)}$, and it is even. So the whole
determinant can be accumulated in Bio~12's orbit order, one $\rho$-orbit at a time, and what has
to be carried across the frontier is: for each open vertex $v$, whether $\bpsi_v$ has already
been spent and whether $\psi_v$ has already been taken. Two bits per vertex, not a set
partition. Section~\ref{sec:wall} explains why this turns out to be the more expensive of the
two alphabets.

\section{The root, and why the box is halved}\label{sec:root}

A reduced determinant needs a root, and $\supp c$ moves with $c$: no vertex of $D_k$ lies in
every support. The diagonal completion \eqref{eq:Aprime} takes care of the vertices that are
missing, and the following lemma says what $\widetilde A^{\,r}_c$ then computes.

\begin{lemma}\label{lem:root}
Let $c$ be a circulation with $c\ne0$. If $r\in\supp c$ then
$\det\widetilde A^{\,r}_c=t(G_c)$, independently of which $r\in\supp c$ is taken; and it
vanishes exactly when $\supp c$ is disconnected. If $r\notin\supp c$ then
$\det\widetilde A^{\,r}_c=0$.
\end{lemma}

\begin{proof}
$c$ is a circulation, so $d^-_c=d^+_c$ and no arc of $G_c$ leaves $\supp c$ or enters it: $A_c$
is block diagonal, $L_c$ on $\supp c$ and the identity off it. Deleting $r\in\supp c$ leaves
$\det$ equal to the corresponding reduced Laplacian of $G_c$, which is $t(G_c)$ by the
matrix-tree theorem; $G_c$ is balanced, so its weak components are strongly connected and the
arborescence count does not depend on the root, and the nullity of $L_c$ on $\supp c$ is the
number of those components \cite[Lemma~3]{Bio12}. Deleting $r\notin\supp c$ leaves the whole
singular block $L_c$ on $\supp c$ intact, so the determinant is $0$.
\end{proof}

So a fixed root is legitimate on any part of the box on which it lies in the support. Such a
part is produced, for free, by the symmetry of the problem.

\begin{lemma}[$\rho$-symmetry]\label{lem:sym}
$c\mapsto m-c$ is an involution of the set of circulations with $c+\rho c=m$, and
$\Nseq(m-c)=\Nseq(c)$.
\end{lemma}

\begin{proof}
$(m-c)(a)=c(\rho a)$, so $a\mapsto\rho a$, $v\mapsto\bar v$ carries $G_{m-c}$ isomorphically to
$G_c$ with all arcs reversed, because $t(\rho a)=\overline{h(a)}$ and $h(\rho a)=\overline{t(a)}$.
Hence $d^+_{m-c}(v)=d^-_c(\bar v)=d^+_c(\bar v)$, so $\Loc$ is unchanged; $\prod_ac(a)!$ is
unchanged because the multiset of values is permuted; and $t(G_{m-c})$ is the number of spanning
\emph{in}-trees of $G_c$, which for a balanced digraph equals the number of out-trees.
\end{proof}

\begin{corollary}\label{cor:half}
Let $j_0$ be a $\rho$-orbit with $m_{j_0}$ odd. Then $f_{j_0}\ne0$ on the box, the involution of
Lemma~\ref{lem:sym} exchanges $\{f_{j_0}>0\}$ with $\{f_{j_0}<0\}$, and
\[
\#\DS(S,k)\;=\;2\!\!\sum_{f:\ f_{j_0}>0}\!\!\Nseq\Big(\frac{m+f}{2}\Big).
\]
On that half-box $c(a_{j_0})\ge1$, so $t(a_{j_0})\in\supp c$ for every point of it and the
single vertex $r=t(a_{j_0})$ serves as root throughout.
\end{corollary}

At $k=10$, $96$ of the $111$ orbits of $\phi$X174 have odd multiplicity, so the corollary
applies; it halves the work and removes the root from the list of things the sweep has to
decide.

\section{What the sweep has to carry besides the fermions}\label{sec:carry}

Two more quantities in \eqref{eq:sum} are not arc-local.

\paragraph{$\Loc$.} $(d^+_c(v)-1)!$ needs the whole out-degree of $v$ at once. The partial
out-degree is therefore carried as well --- but only on the sub-frontier of the vertices whose
out-arcs are not all decided, which at $k=10$ is $18$ wide against the $26$ of the frontier
itself, and is narrower still under an order chosen with it in mind. It is spent, and dropped,
at the step at which the last out-arc of $v$ is decided; if that out-degree is $0$ the vertex is
outside $\supp c$ and the factor $(1-\bpsi_v\psi_v)$ of \eqref{eq:prod} is applied instead.

\paragraph{Integrality.} $\prod_ac(a)!$ has no reason to divide anything along the way. It is
removed from the sweep altogether: on a free orbit $c(a_j)+c(\rho a_j)=m_j$, so
\begin{equation}\label{eq:int}
\prod_{a}\frac1{c(a)!}\;=\;\frac{\prod_j\binom{m_j}{c(a_j)}}{\prod_j m_j!\,\cdot
\prod_{\rho a=a}(m_a/2)!},
\end{equation}
whose numerator is a product of integer weights, one per orbit, and whose denominator is a
constant of $D_k$. The sweep therefore runs over $\Z$ and divides once, at the end.

\section{The sweep}\label{sec:sweep}

The order of the orbits is Bio~12's greedy order, re-tuned: the elements the frontier now has
to carry are one $\psi$ per vertex over its whole window, weighted for the imbalance it travels
with, and one $\bpsi$ per vertex over its out-arc window, weighted for the partial out-degree.
A short local search on top of a randomised greedy brings $\phi$X174 at $k=10$ from $42$
elements to $35$.

The state of the sweep is
\[
\big(\text{imbalance }b_v\big)_{v\ \text{open}}\ \times\
\big(\text{partial out-degree}\big)_{v\ \text{out-open}}\ \times\
\big(\text{which }\bpsi_v,\psi_v\ \text{are spent}\big),
\]
and its value is an integer. Each vertex is given a \emph{slot} for the length of its window, by
colouring the windows as intervals; slots are reused, so the layout is fixed once and for all
and nothing has to be relabelled from step to step. The imbalances live in fixed-width fields of
one machine word, biased so that an empty field reads zero, and the capacity bound of Bio~12
---
$|b_v|$ at most the capacity still to come --- is tested on all fields at once by two additions
and two masks, which also forces a vertex leaving the frontier to balance. The out-arc window of
$v$ sits inside its full window, so $\bpsi_v$ and $\psi_v$ can share $v$'s slot; giving them the
two bits $2s$ and $2s+1$ of that slot makes the generator order interleaved, which by
Remark~\ref{rem:interleave} is what the constant of Proposition~\ref{prop:arc} requires. The
signs are then local: multiplying on the right by a generator costs $(-1)$ to the number of
spent generators above it, and integrating one out when its window closes costs $(-1)$ to the
number of spent generators below it, both read off the bit mask; the permutation between the
order in which the generators leave and the canonical one is a single sign computed once.

Finally, the half-box may be cut into disjoint sub-boxes by pinning the values of a few early
orbits, and the sums added. This changes nothing but the peak memory.

The sweep reproduces $\#\DS=10$ at $k=12$ and $864$ at $k=11$ from the half-box, with one
fixed root and one division at the end. It also agrees with the brute-force sum on twelve
random covers. At $k=11$ it peaks at $560$ states.

\section{Why the frontier stops short of $k=10$}\label{sec:wall}

The frontier state of Section~\ref{sec:sweep} is \emph{linear}. At a fixed imbalance it is a
vector in the exterior algebra of the open generators, and summing over histories adds vectors.
The alphabet is small, but the vectors are not sparse. A boundary vertex $v$ whose out-window
is still open, and which already carries an arc to an open vertex $u$, contributes the factor
$1-\bpsi_v\psi_v+\bpsi_v\psi_u$. A single history with $s$ such vertices therefore already has
up to $3^{s}$ non-zero coordinates, before any two histories are added. At $k=10$ the tuned
order still has $35$ open generators at its widest, and the weighted width estimate of the
order search is $2^{86.8}$. The cancellations of the Berezin integral occur only at the end;
along the way, the linear state pays for all the signed terms that will cancel.

The obvious alternative is the positive expansion of the same determinant,
$t(G_c)=\sum_T\prod_{a\in T}c(a)$ over in-arborescences $T$ of $\supp c$ rooted at $r$. It needs
no signs, and it puts \emph{one} state on each history.

\begin{proposition}[arborescence elimination]\label{prop:arbor}
Process the orbits in order. Call a vertex \emph{pending} until its last out-arc has been
decided, and let every pending vertex be the root of its own class. Carry, for each open vertex,
its imbalance. For a pending vertex, also carry its candidate arcs so far as a multiset over the
classes their heads belong to; its partial out-degree is their total. For a vertex that is no
longer pending, carry the class it belongs to. When the last out-arc of a pending vertex
$v\ne r$ has been decided, do one of two things. If its out-degree is $0$, declare it outside
$\supp c$. Otherwise let it choose a candidate class $K$ other than its own, with weight the
multiplicity towards $K$ times $(d^+_c(v)-1)!$, and merge its class into $K$. The root never
chooses; it contributes $(d^+_c(r)-1)!$. Then the sum of the weights over all runs equals
$\sum_c\prod_j\binom{m_j}{c(a_j)}\,t(G_c)\Loc(G_c)$ over the half-box.
\end{proposition}

\begin{proof}
A choice is an arc from $v$ into a class other than $v$'s own. Following the choices from any
vertex eventually reaches a pending vertex or $r$, so a cycle in the chosen arcs would have a
last vertex to choose; that vertex would have found its head already in its own class.
Conversely, a choice into a different class never closes a cycle. So the runs that survive
correspond exactly to the in-arborescences of $\supp c$ to $r$, each counted with
$\prod_{a\in T}c(a)$. A component of the support that does not contain $r$ has a last pending
vertex whose candidates all lie in its own class, and so it contributes nothing.
\end{proof}

It is exact. It reproduces $10$ and $864$, the random covers, and itself as a sum over
sub-boxes. At $k=10$ it agrees with the point-by-point sum of $\Nseq$ on five random sub-boxes of
the half-box. But it merges too few histories. Those five sub-boxes hold between $12$ and $166$
points, with $30$ to $36$ of the $111$ orbits pinned, and they already peak at between $55\,890$
and $623\,000$ states. The candidate multisets and class pointers are strictly finer than the
connectivity partition with which Bio~12 counted the connected points, and that elimination
already needed $9.7$ million states. Neither frontier reaches the whole half-box on a
$16$~GB machine.

\section{Point by point}\label{sec:points}

What does reach it is the least clever method, made cheap at each point. The half-box has
$1\,240\,843\,950$ points: half of Bio~12's count, by Corollary~\ref{cor:half}. The support of a
point has about $126$ arcs on about $110$ vertices, so most vertices have out-degree $1$, and
the determinant is almost entirely forced.

\begin{lemma}[forced arcs]\label{lem:forced}
Let $G$ be a balanced multidigraph, $r$ a vertex, and $v\ne r$ a vertex all of whose out-arcs
(there are $w\ge1$ of them) end at one vertex $u\ne v$. Then $t_r(G)=w\,t_r(G/vu)$, where $G/vu$
identifies $v$ with $u$, deletes the $w$ arcs $v\to u$ and redirects the arcs into $v$ to $u$.
\end{lemma}

\begin{proof}
Every in-arborescence to $r$ contains exactly one out-arc of $v$, hence one of the $w$ arcs to
$u$. Contracting it is a bijection onto the in-arborescences of $G/vu$ to $r$; the other copies
of $v\to u$, and the arcs from $u$ to $v$, become loops, which no arborescence uses.
\end{proof}

Applying the lemma until no vertex qualifies leaves a reduced Laplacian of small order. At each
point, then:
\begin{enumerate}
\item the support is tested for connectivity by union-find, and discarded if it is not
connected (Lemma~\ref{lem:root});
\item $t(G_c)$ is computed by Lemma~\ref{lem:forced} followed by Gaussian elimination on
what is left;
\item $\Loc(G_c)/\prod_a c(a)!$ is read from tables.
\end{enumerate}
All of this is done modulo the five largest primes below $2^{31}$, so every product fits in a
machine word. The points are visited by Bio~8's imbalance-bounded backtracking in Bio~12's
greedy order, compiled with \texttt{numba}. The first $35$ levels are expanded into $28\,512$
prefixes, which are shuffled and processed in parallel. The half-sum is recovered by the
Chinese remainder theorem. This is exact: the half-sum is at most half of Bio~8's upper bound
$1.368\cdot10^{36}$, far below the product of the five primes (about $2^{155}$).

The run is checked in three ways, all exactly:
\begin{enumerate}
\item At $k=12$ and $k=11$ it returns $5$ and $264$ half-box points, $4$ and $109$ of them
connected, and $\#\DS=10$ and $864$; on twelve random covers it agrees with the brute force.
\item On $k=10$ sub-boxes it agrees with the arborescence elimination of
Proposition~\ref{prop:arbor}. That is a second, unrelated algorithm, and the sub-boxes are
built from the same order, windows and root as a full run would use.
\item At $k=10$ it counts exactly half of Bio~12's box points and exactly half of its
connected points (Corollary~\ref{cor:half}; the involution preserves connectivity and has no
fixed point on the box).
\end{enumerate}

\section{$\phi$X174 at $k=10$}\label{sec:result}

The root is vertex $20$, the tail of orbit $33$, which has multiplicity $1$. The enumeration
visits
\[
1\,240\,843\,950=\tfrac12\cdot 2\,481\,687\,900 \ \text{points, of which}\
454\,551\,605=\tfrac12\cdot 909\,103\,210\ \text{have connected support},
\]
with no overflow of the work arrays. It takes about thirty-five minutes on sixteen threads.
The sum of $\Nseq$ over the half-box, doubled, is
\[
\boxed{\;\#\DS(\phi\mathrm{X}174,10)=31\,925\,753\,246\,212\approx3.19\cdot10^{13}.\;}
\]
This lies inside Bio~8's bounds $8\,610\,708\,632\le\#\DS\le1.368\cdot10^{36}$: about $3\,700$
times the lower bound and twenty-three orders of magnitude below the upper one. Averaged over
the $909\,103\,210$ points with connected support, a point carries about $35\,118$ circular
sequences. At $k=11$ the average is $864/218\approx4$.

So at word length $10$, $\phi$X174 is one of about $3.2\cdot10^{13}$ circular sequences with
the same double-stranded $10$-mer spectrum. The lattice has rank $48$, but the count is not
astronomical: the box has $2.5\cdot10^9$ points, and each carries a modest number of Eulerian
circuits.

\begin{table}[t]
\centering
\begin{tabular}{crrrr}
\toprule
$k$ & $\rank\Lm$ & box points & connected & $\#\DS$ \\
\midrule
12 & 4  & 10 & 8 & 10 \\
11 & 12 & 528 & 218 & 864 \\
10 & 48 & 2\,481\,687\,900 & 909\,103\,210 & 31\,925\,753\,246\,212 \\
\bottomrule
\end{tabular}
\caption{Bio~8's Table~2, completed. The first two columns at $k=10$ are Bio~12's; the last is
this note's.}\label{tab:done}
\end{table}

\section*{Ancillary files}
\texttt{ds\_fermion.py} is standalone: Python 3.10+ and the standard library. It rebuilds the
double cover from its definition. It checks Proposition~\ref{prop:arc} against the reduced
determinant point by point on $\phi$X174 at $k=12$ and on twelve small random covers, and it
checks Lemma~\ref{lem:root}, Lemma~\ref{lem:sym} and Corollary~\ref{cor:half} on the enumerable
boxes. It reproduces $\#\DS=10$ and $864$ with the fermionic sweep and with the arborescence
elimination, and it checks the latter, and Lemma~\ref{lem:forced}, on $k=10$ sub-boxes: seven
checks in about ten minutes.

\texttt{ds\_count.py} needs \texttt{numpy} and \texttt{numba}, and imports
\texttt{ds\_fermion.py}. It is the enumeration of Section~\ref{sec:points} with its three
checks, and it computes $\#\DS(\phi\mathrm{X}174,10)$. The run takes about thirty-five minutes on
sixteen threads and resumes from a checkpoint.

\begin{thebibliography}{9}
\bibitem{Bio8} R.~Chen, \emph{Double-stranded reconstruction and the lattice of a
reverse-complement double cover} (Bio~8 / XXI).
\bibitem{Bio12} R.~Chen, \emph{The double-stranded lattice sum by frontier elimination}
(Bio~12 / XXV).
\end{thebibliography}

\end{document}
